% ============================================================================
% Bab 6: Pernyataan Matematis Formal
% Sumber buku: part1-linear-programming/ch05-lp-theory/formalize-LP.tex
% Penomoran latihan diselaraskan dengan urutan \begin{ex} per 2026-07-13.
% Pemeriksaan numerik diverifikasi dengan numpy/scipy.
% ============================================================================
\section{Bab 6: Pernyataan Matematis Formal}

\textbf{Penilaian cakupan.} Ke-12 latihan memberikan tinjauan yang seimbang atas bab ini: latihan pemanasan (6.1--6.5) melatih vektor, kombinasi, rank, pengenalan konveksitas, dan bentuk matriks suatu program linier; soal inti (6.6--6.9) membahas secara konkret kombinasi konveks, setengah ruang, serta perangkat titik sudut/titik ekstrem pada polihedron yang sama dengan yang digunakan untuk mengilustrasikan Teorema Representasi, suatu pilihan yang bagus; kelompok konsep (6.10--6.11) menguji pembuktian (gabungan versus irisan, dan tepat di mana keterbatasan digunakan dalam teorema titik sudut); dan satu-satunya soal tantangan (6.12) merupakan latihan pembuktian yang sesungguhnya. Cakupan pembuktian sengaja dibuat ringan, sesuai dengan kotak peringatan dalam bab ini bahwa mahasiswa tidak perlu menghasilkan kembali pembuktian-pembuktian tersebut; hal ini sesuai dan penulis telah menyatakan bahwa pendekatan ini dapat diterima untuk bab-bab teori. Tidak ditemukan kesalahan data; keempat Solusi Pilihan dalam buku (6.1, 6.6, 6.9, 6.10) telah diperiksa dan digunakan kembali di bawah ini. Tidak ada latihan di sini yang tampaknya berasal dari sumber luar; isi dan latihan tampak asli dari buku ini. Satu catatan pedagogis kecil: pada Latihan 6.2 bagian 1, setiap vektor dalam $\mathbb{R}^2$ merupakan kombinasi linier dari $\mathbf{v}^1, \mathbf{v}^2$, sehingga ketiga jawabannya adalah “ya”; hal ini dapat dianggap sebagai inti latihan tersebut, jadi tidak ditandai sebagai masalah.

\exsol{6.1}{Operasi Vektor}{1}
Dengan menghitung per komponen (sesuai dengan Solusi Pilihan dalam buku):
\begin{enumerate}
\item $\mathbf{v} + \mathbf{u} = [3 + (-1),\; -2 + 0,\; 4 + 2] = [2, -2, 6]$.
\item $2\mathbf{v} = [6, -4, 8]$.
\item $\mathbf{v} \cdot \mathbf{u} = (3)(-1) + (-2)(0) + (4)(2) = -3 + 0 + 8 = 5$.
\end{enumerate}

\exsol{6.2}{Kombinasi Linier dan Konveks}{1}
\begin{enumerate}
\item Ketiganya merupakan kombinasi linier dari $\mathbf{v}^1 = [1,0]$ dan $\mathbf{v}^2 = [0,1]$, karena kedua vektor ini merentang $\mathbb{R}^2$: sembarang $[a, b] = a\,\mathbf{v}^1 + b\,\mathbf{v}^2$. Secara eksplisit,
\begin{itemize}
\item[(a)] $[2,3] = 2\,\mathbf{v}^1 + 3\,\mathbf{v}^2$,
\item[(b)] $[1,1] = 1\,\mathbf{v}^1 + 1\,\mathbf{v}^2$,
\item[(c)] $[1,-1] = 1\,\mathbf{v}^1 - 1\,\mathbf{v}^2$.
\end{itemize}
Perhatikan bahwa (c) adalah kombinasi linier, tetapi bukan kombinasi nonnegatif.
\item Ya. Gunakan pasangan $\mathbf{v}^1, \mathbf{v}^2$ dengan $\lambda = \tfrac12$:
\[
\tfrac12 [1, 0] + \tfrac12 [0, 1] = [\tfrac12, \tfrac12],
\]
dan bobotnya nonnegatif serta berjumlah $1$. (Ini satu-satunya pasangan yang berhasil: kombinasi dari $\mathbf{v}^1$ dan $\mathbf{v}^3$ memiliki koordinat pertama $1$, sedangkan kombinasi dari $\mathbf{v}^2$ dan $\mathbf{v}^3$ memiliki koordinat kedua $1$.)
\end{enumerate}

\exsol{6.3}{Kebebasan Linier dan Rank Matriks}{1}
\begin{enumerate}
\item Tidak. Baris kedua $[2,4]$ sama dengan $2$ kali baris pertama $[1,2]$, sehingga baris-baris tersebut bergantung linier: $2[1,2] - [2,4] = [0,0]$ adalah kombinasi taktrivial yang bernilai nol.
\item Tidak. Kolom kedua $[2,4]^\top$ sama dengan $2$ kali kolom pertama $[1,2]^\top$, sehingga kolom-kolom tersebut juga bergantung.
\item $\operatorname{rank}(A) = 1$: jumlah maksimum baris yang bebas linier (atau, secara ekuivalen, kolom yang bebas linier) adalah $1$. Pemeriksaan yang konsisten: $\det A = 1\cdot 4 - 2 \cdot 2 = 0$, sehingga rank kurang dari $2$, dan $A \neq 0$, sehingga rank sekurang-kurangnya $1$.
\end{enumerate}

\exsol{6.4}{Konveks atau Tidak?}{1}
\begin{enumerate}
\item \emph{Cakram $D$: konveks.} Jika $\x, \y \in D$, maka berdasarkan pertaksamaan segitiga untuk norma Euklides, $\|\lambda \x + (1-\lambda)\y\| \leq \lambda\|\x\| + (1-\lambda)\|\y\| \leq \lambda\cdot 2 + (1-\lambda)\cdot 2 = 2$, sehingga ruas garisnya tetap berada dalam $D$. (Bola bersifat konveks, sebagaimana tercantum dalam bab ini.)
\item \emph{Cincin $R$: tidak konveks.} Ambil $\x = (1, 0)$ dan $\y = (-1, 0)$; keduanya memiliki norma $1$, sehingga keduanya berada dalam $R$. Titik tengahnya adalah $(0,0)$, dengan $x_1^2 + x_2^2 = 0 < 1$, yang tidak berada dalam $R$: ruas garis tersebut melewati lubang.
\item \emph{Setengah ruang $H$: konveks.} Untuk $\x, \y$ dengan $x_1 + 2x_2 \leq 3$ dan $y_1 + 2y_2 \leq 3$, linearitas memberikan $(\lambda \x + (1-\lambda)\y)_1 + 2(\lambda \x + (1-\lambda)\y)_2 = \lambda(x_1 + 2x_2) + (1-\lambda)(y_1 + 2y_2) \leq \lambda\cdot 3 + (1-\lambda)\cdot 3 = 3$. (Setengah ruang bersifat konveks, sebagaimana tercantum dalam bab ini.)
\item \emph{Himpunan $S = \{\x : |x_1| \geq 1\}$: tidak konveks.} Ambil $\x = (1, 0) \in S$ dan $\y = (-1, 0) \in S$. Titik tengahnya $(0,0)$ memiliki $|x_1| = 0 < 1$, sehingga tidak berada dalam $S$: himpunan tersebut terdiri atas dua setengah bidang yang saling lepas dan ruas garisnya melintasi celah di antara keduanya.
\end{enumerate}

\exsol{6.5}{Menulis Program Linier dalam Notasi Matriks}{1}
Urutkan kendala sebagai berikut: baris lemon, baris air, lalu dua kendala tanda yang ditulis ulang, yaitu $-x_1 \leq 0$ dan $-x_2 \leq 0$. Dengan demikian,
\[
\c = \begin{bmatrix} 3 \\ 2 \end{bmatrix}, \qquad
A = \begin{bmatrix} 1 & 3 \\ 2 & 1 \\ -1 & 0 \\ 0 & -1 \end{bmatrix}, \qquad
\b = \begin{bmatrix} 6 \\ 4 \\ 0 \\ 0 \end{bmatrix},
\]
dan program liniernya adalah $\max\{ \c^\top \x : A\x \leq \b \}$.

Verifikasi pada $\x = (1.2,\, 1.6)$:
\[
A\x = \begin{bmatrix} 1.2 + 3(1.6) \\ 2(1.2) + 1.6 \\ -1.2 \\ -1.6 \end{bmatrix}
= \begin{bmatrix} 6 \\ 4 \\ -1.2 \\ -1.6 \end{bmatrix} \leq \begin{bmatrix} 6 \\ 4 \\ 0 \\ 0 \end{bmatrix},
\]
sehingga $\x$ layak (diverifikasi secara numerik). Dua kendala pertama berlaku dengan kesamaan: keduanya adalah kendala aktif pada titik sudut ini, yang merupakan solusi optimal dari program linier Penjual Limun pada bab sebelumnya.

\exsol{6.6}{Kombinasi Konveks dan Geometri}{2}
Segitiga $P$ adalah himpunan kombinasi konveks dari $(0,0)$, $(2,0)$, $(0,2)$; secara ekuivalen, $P = \{ \x \in \mathbb{R}^2 : x_1 \geq 0,\ x_2 \geq 0,\ x_1 + x_2 \leq 2 \}$. (Sesuai dengan Solusi Pilihan dalam buku.)
\begin{enumerate}
\item Gambarnya adalah segitiga tertutup dengan ketiga titik sudut yang diberikan: kedua kakinya terletak pada sumbu-sumbu koordinat dan hipotenusanya adalah ruas dari $x_1 + x_2 = 2$ yang menghubungkan $(2,0)$ dengan $(0,2)$.
\item Ya. Ambil $\lambda_1 = 0$, $\lambda_2 = \lambda_3 = \tfrac12$:
\[
0 \begin{bmatrix} 0 \\ 0 \end{bmatrix} + \tfrac12 \begin{bmatrix} 2 \\ 0 \end{bmatrix} + \tfrac12 \begin{bmatrix} 0 \\ 2 \end{bmatrix} = \begin{bmatrix} 1 \\ 1 \end{bmatrix},
\]
dengan bobot nonnegatif yang berjumlah $1$; jadi, $(1,1)$ adalah titik tengah hipotenusa.
\item Tidak. Setiap kombinasi titik-titik sudut yang menghasilkan $(1.5, 1.5)$ harus memenuhi $2\lambda_2 = 1.5$ dan $2\lambda_3 = 1.5$, sehingga $\lambda_2 = \lambda_3 = 0.75$ dan $\lambda_2 + \lambda_3 = 1.5 > 1$. Secara ekuivalen, $1.5 + 1.5 = 3 > 2$ melanggar $x_1 + x_2 \leq 2$, sehingga $(1.5, 1.5) \notin P$.
\end{enumerate}

\exsol{6.7}{Geometri dan Konveksitas Setengah Ruang}{2}
\begin{enumerate}
\item Batasnya adalah garis $x_1 + 2x_2 = 4$, yang melalui $(4, 0)$ dan $(0, 2)$. Setengah ruang $H$ adalah garis tersebut beserta semua titik di bawahnya (sisi yang memuat titik asal, karena $0 + 0 \leq 4$).
\item $(2,1)$: $2 + 2(1) = 4 \leq 4$, jadi ya, $(2,1) \in H$; titik tersebut tepat berada pada batas. $(1,2)$: $1 + 2(2) = 5 > 4$, sehingga $(1,2) \notin H$.
\item Titik tengah dari $(0,0)$ dan $(2,1)$ adalah $(1, \tfrac12)$, dan $1 + 2\cdot\tfrac12 = 2 \leq 4$, sehingga titik tengah tersebut berada dalam $H$, konsisten dengan konveksitas. (Secara ketat, memeriksa satu titik tengah mengilustrasikan konveksitas tetapi tidak membuktikannya; argumen umumnya adalah perhitungan linearitas dalam Teorema “Konveksitas Polihedron” yang diterapkan pada baris tunggal $[1\;\; 2]$: untuk sembarang $\x, \y \in H$ dan $\lambda \in [0,1]$, $(\lambda\x + (1-\lambda)\y)$ memenuhi $x_1 + 2x_2 \leq \lambda 4 + (1-\lambda)4 = 4$.)
\end{enumerate}

\exsol{6.8}{Memeriksa Titik Ekstrem}{2}
Tuliskan kelima kendala dengan baris $A_1 = [1, -4]$, $A_2 = [-1, 1]$, $A_3 = [-1, 2]$, $A_4 = [-1, 0]$, $A_5 = [0, -1]$ dan ruas kanan $(0, 1, 4, 0, 0)$.
\begin{enumerate}
\item Pada $(2,3)$: $A_1\x = 2 - 12 = -10 \leq 0$; $A_2\x = -2 + 3 = 1$ (aktif); $A_3\x = -2 + 6 = 4$ (aktif); $-x_1 = -2 \leq 0$; $-x_2 = -3 \leq 0$. Jadi, $(2,3) \in P$ dan kendala yang aktif tepat baris $2$ dan $3$. Matriks baris-baris aktif,
\[
A_I = \begin{bmatrix} -1 & 1 \\ -1 & 2 \end{bmatrix}, \qquad \det A_I = (-1)(2) - (1)(-1) = -1 \neq 0,
\]
memiliki rank $2$ (diverifikasi secara numerik), sehingga kedua baris aktif itu bebas linier dan $|I| = n = 2$. Sistem $-x_1 + x_2 = 1$, $-x_1 + 2x_2 = 4$ memiliki solusi tunggal $(2,3)$, jadi $(2,3)$ adalah titik sudut $P$, dan karenanya merupakan titik ekstrem berdasarkan teorema ekuivalensi dalam bab ini.
\item Pada $(1,2)$: $A_1\x = 1 - 8 = -7 \leq 0$; $A_2\x = -1 + 2 = 1$ (aktif); $A_3\x = -1 + 4 = 3 < 4$; $-1 < 0$; $-2 < 0$. Jadi, $(1,2) \in P$ dan hanya kendala $2$ yang aktif; satu baris tidak dapat memuat $2$ baris yang bebas linier, sehingga $(1,2)$ bukan titik sudut. Untuk menyatakannya sebagai kombinasi konveks, bergeraklah sepanjang garis aktif $-x_1 + x_2 = 1$: titik $(0,1)$ dan $(2,3)$ keduanya berada dalam $P$ (periksa $(0,1)$: $0 - 4 = -4 \leq 0$, $-0+1 = 1 \leq 1$, $-0+2 = 2 \leq 4$, $0 \geq 0$, $1 \geq 0$), dan
\[
(1, 2) = \tfrac12 (0, 1) + \tfrac12 (2, 3),
\]
yang merupakan kombinasi konveks ketat dari dua titik lain dalam $P$. Dengan demikian, $(1,2)$ bukan titik ekstrem.
\end{enumerate}

\exsol{6.9}{Kombinasi Konveks Titik-Titik Ekstrem}{2}
(Sesuai dengan Solusi Pilihan dalam buku.) Koordinat pertama mengharuskan $2\lambda_3 = 1$, sehingga $\lambda_3 = \tfrac12$. Koordinat kedua kemudian memberikan $\lambda_2 + \tfrac32 = \tfrac74$, sehingga $\lambda_2 = \tfrac14$, dan normalisasi memberikan $\lambda_1 = 1 - \tfrac14 - \tfrac12 = \tfrac14$. Semua pengali nonnegatif, dan memang
\[
\tfrac14 \begin{bmatrix} 0 \\ 0 \end{bmatrix} + \tfrac14 \begin{bmatrix} 0 \\ 1 \end{bmatrix} + \tfrac12 \begin{bmatrix} 2 \\ 3 \end{bmatrix} = \begin{bmatrix} 1 \\ 7/4 \end{bmatrix}
\]
(diverifikasi secara numerik). Tidak diperlukan arah ekstrem karena $(1, \tfrac74)$ sudah berada dalam selubung konveks dari ketiga titik ekstrem; dalam Teorema Representasi kita dapat mengambil semua $\mu_j = 0$. Arah ekstrem baru diperlukan untuk titik-titik $P$ di luar selubung konveks tersebut, yang terletak lebih jauh pada bagian tak berbatas dari daerah itu.

\exsol{6.10}{Apakah Gabungan Himpunan Konveks Bersifat Konveks?}{2}
(Sesuai dengan Solusi Pilihan dalam buku.)
\begin{enumerate}
\item Ambil $C_1$ sebagai cakram satuan yang berpusat di $(-2, 0)$ dan $C_2$ sebagai cakram satuan yang berpusat di $(2, 0)$; masing-masing cakram bersifat konveks. Titik $(-2, 0) \in C_1$ dan $(2, 0) \in C_2$ berada dalam gabungannya, tetapi titik tengahnya $(0,0)$ berjarak $2 > 1$ dari kedua pusat, sehingga tidak berada dalam kedua cakram tersebut. Oleh karena itu, $C_1 \cup C_2$ tidak konveks.
\item Jika satu himpunan memuat himpunan yang lain, misalnya $C_1 \subseteq C_2$, maka $C_1 \cup C_2 = C_2$, yang bersifat konveks. Keterangkuman ini cukup, tetapi tidak perlu: dengan $C_1 = \{\x : x_2 \geq 0\}$ dan $C_2 = \{\x : x_2 \leq 0\}$, tidak satu pun dari kedua setengah bidang memuat yang lain, tetapi $C_1 \cup C_2 = \mathbb{R}^2$ bersifat konveks.
\end{enumerate}

\exsol{6.11}{Letak Penggunaan Keterbatasan dalam Teorema Titik Sudut}{2}
\begin{enumerate}
\item Keterbatasan digunakan tepat satu kali, setelah pembuktian menetapkan $\c^\top \mathbf{d} = 0$: pembuktian menyatakan bahwa “karena $P$ terbatas, garis $\{\x(t) : t \in \mathbb{R}\}$ tidak mungkin seluruhnya termuat dalam $P$, sehingga himpunan $\{t \geq 0 : \x(t) \in P\}$ adalah interval tertutup dan terbatas $[0, t^+]$,” dan pada $t = t^+$ suatu kendala baru menjadi aktif. Tanpa keterbatasan, seluruh garis $\x^* + t\mathbf{d}$ dapat tetap layak untuk semua $t$; dalam hal itu, tidak ada kendala baru yang pernah menjadi aktif, jumlah kendala aktif yang bebas linier tidak pernah bertambah, dan induksi terhenti. (Lebih buruk lagi, polihedron mungkin sama sekali tidak memiliki titik sudut, seperti pada bagian 2.)
\item Misalkan $P = \{\x \in \mathbb{R}^2 : 0 \leq x_2 \leq 1\}$, sebuah lempeng horizontal, dan maksimalkan $-x_2$.

\emph{Tidak ada titik sudut.} Baris kendalanya adalah $[0, 1]$ (untuk $x_2 \leq 1$) dan $[0, -1]$ (untuk $-x_2 \leq 0$). Pada sembarang titik dalam $P$, paling banyak satu di antaranya yang dapat aktif ($x_2$ tidak dapat sama dengan $0$ sekaligus $1$), dan bahkan keduanya bersama-sama hanya merentang arah sumbu $x_2$. Jadi, kendala-kendala aktif pada sembarang titik memiliki rank paling tinggi $1 < 2 = n$, dan tidak ada titik dalam $P$ yang merupakan solusi tunggal dari $n$ kendala aktif yang bebas linier: $P$ tidak memiliki titik sudut. (Secara geometris, $P$ memuat garis horizontal yang melalui setiap titiknya, dan himpunan yang memuat sebuah garis tidak memiliki titik ekstrem.)

\emph{Optimum ada.} Fungsi tujuan $-x_2$ paling besar bernilai $0$ pada $P$, dan sama dengan $0$ tepat ketika $x_2 = 0$. Jadi, program linier tersebut memiliki nilai optimal $0$ dan himpunan optimumnya adalah seluruh garis batas $\{(x_1, 0) : x_1 \in \mathbb{R}\}$, yaitu tak terhingga banyak optimum, tetapi tidak satu pun pada titik sudut.

\emph{Hipotesis yang gagal.} $P$ tak berbatas ($x_1$ bebas), sehingga hipotesis keterbatasan dalam teorema tidak terpenuhi dan teorema tersebut memang tidak berlaku; tidak ada kontradiksi. Contoh ini menunjukkan bahwa hipotesis tersebut tidak dapat dihapus: daerah layak yang tak berbatas mungkin tidak memiliki titik sudut tempat optimum dapat dicapai, bahkan ketika optimum ada.
\end{enumerate}

\exsol{6.12}{Irisan Keluarga Himpunan Konveks Sembarang}{3}
\emph{Bukti.} Misalkan $\x, \y \in C = \bigcap_{i \in I} C_i$ dan $\lambda \in [0, 1]$; tetapkan $\mathbf{z} = \lambda\x + (1 - \lambda)\y$. Ambil sembarang indeks $i \in I$. Karena $\x \in C$, kita mempunyai $\x \in C_i$, dan demikian pula $\y \in C_i$. Karena $C_i$ konveks, $\mathbf{z} \in C_i$. Karena $i \in I$ dipilih sembarang, $\mathbf{z}$ termasuk dalam setiap $C_i$, yaitu $\mathbf{z} \in C$. Dengan demikian, $C$ memuat semua kombinasi konveks dari titik-titiknya dan bersifat konveks. Tidak ada bagian argumen yang menggunakan batas apa pun pada $|I|$; argumen ini berlaku persis sama untuk keluarga tak hingga. Jika $C$ kosong, pernyataan tersebut berlaku secara vakum karena tidak ada pasangan $\x, \y$ yang perlu diuji. $\blacksquare$

\emph{Penerapan pada polihedron.} Polihedron $P = \{\x \in \mathbb{R}^n : A\x \leq \b\}$ dengan $m$ baris tepat merupakan irisan dari $m$ setengah ruang $H_i = \{\x : A_i \x \leq b_i\}$. Setiap setengah ruang bersifat konveks (kasus satu baris dari argumen linearitas, atau daftar himpunan konveks dalam bab ini). Berdasarkan hasil yang baru saja dibuktikan, yang diterapkan pada keluarga berhingga $\{H_1, \dots, H_m\}$, irisan $P$ bersifat konveks. Ini membuktikan kembali Teorema “Konveksitas Polihedron” tanpa mengulangi perhitungan per komponen.
